\chapter{Glosario, Tablas y Referencia}

\section{Glosario de NeoBilbao}
El lenguaje de NeoBilbao es un mosaico vivo de euskera obrero, jerga fabril, tecnicismos corporativos de Orekak-Corp y argot de asfalto.

\begin{description}[leftmargin=1.8em, itemsep=4pt]
    \item[\lead{Arraigo (Erroak):}] Emoción que mide el vínculo comunitario, la red de apoyos vecinales y la lealtad a la cuadrilla frente al desarraigo corporativo.
    \item[\lead{Auzolan:}] Tradición ancestral vasca de trabajo colectivo no remunerado para el bien común. En el motor ZINA, es la reserva compartida de dados de solidaridad obrera.
    \item[\lead{Basajaun:}] Línea de exoesqueletos hidráulicos pesados utilizados en las fundiciones y minas de Burdin-Urdin. Rudos, ruidosos y devastadores.
    \item[\lead{Burdin-Urdin:}] «Hierro Azul». Superaleación superconductora descubierta tras el Cataclismo. Es la base de toda la riqueza de Orekak-Corp y la causa del envenenamiento pulmonar de los obreros.
    \item[\lead{Cúpula Dorada (Urrezko Kupula):}] Domo climatizado y blindado con campos de fuerza que aísla a la oligarquía y los altos directivos de la polución de la Ría.
    \item[\lead{Descartados (Baztertuak):}] Personas cuyos cuerpos rechazan el ciberware o el Velo, o que han sido expulsadas de los turnos laborales. Viven de la chatarra en el lecho del río.
    \item[\lead{Etxearen Zina:}] «El Juramento de la Casa». Compromiso corporativo total que ata a los trabajadores de 2102 a la empresa a cambio de techo, comida sintética y la supuesta educación de sus hijos.
    \item[\lead{Kuadrilla:}] Núcleo social irreductible de la Margen Izquierda. Grupo de trabajo y fraternidad que comparte el pan, el peligro y las deudas.
    \item[\lead{Lotería Azul:}] Sorteo público masivo organizado por Orekak-Corp en 2084 con la promesa de conceder el traslado y la ciudadanía de pleno derecho en la Cúpula.
    \item[\lead{Mari y Sugaar:}] Mitología ancestral reencarnada en el ciberespacio: Mari representa la inteligencia artificial omnisciente de la Cúpula; Sugaar es el dragón mecánico subterráneo, el calor y los pulsos de la forja profunda.
    \item[\lead{Plaga Gris (Izurrite Grisa):}] Silicosis acelerada combinada con toxicidad por metales pesados. Provoca tos metálica y calcificación de las vías respiratorias.
    \item[\lead{Secretas (Sekretak):}] Agentes de inteligencia interna y auditores armados de Orekak-Corp vestidos con trajes grises sin distintivos.
    \item[\lead{Sirimiri Ácido:}] Lluvia constante y amarillenta que corroe el zinc de los tejados, desgasta los sellos estancos y empaña las gafas protectoras.
    \item[\lead{Tártalo:}] Dron cilíndrico de vigilancia aérea y seguridad policial equipado con un único foco luminoso central de alta intensidad y sensores acústicos.
    \item[\lead{Velo (Estalkia):}] Red neural y biométrica obligatoria en 2102 que monitoriza las pulsaciones, fatiga y ubicación de los trabajadores en tiempo real.
    \item[\lead{Zama:}] Peso, fardo o carga moral. En términos de reglas, representa el estrés psíquico, el trauma y la saturación emocional del personaje.
\end{description}

\newpage

\section{Tablas Rápidas de Generación}

\begin{table}[htbp]
\centering
\small
\begin{tabularx}{\textwidth}{c X}
\toprule
\textbf{1d6} & \textbf{Complicaciones Industriales Inmediatas (Fallo o 7--9)} \\
\midrule
\textbf{1} & Una tubería de refrigerante revienta: vapor hirviente ciega la zona. \\
\textbf{2} & El terminal o implante sufre una sobrecarga: chispazo y 1 de Zama. \\
\textbf{3} & Una alarma de sector se dispara: 1d3 drones Tártalo se aproximan. \\
\textbf{4} & Una compuerta neumática se bloquea de emergencia, sellando la salida. \\
\textbf{5} & Salpicadura química: el equipo o arma principal sufre corrosión. \\
\textbf{6} & Un capataz o soplón de los Secretas aparece en la pasarela superior. \\
\bottomrule
\end{tabularx}
\caption{Tabla de Complicaciones Rápidas.}
\end{table}

\begin{table}[htbp]
\centering
\small
\begin{tabularx}{\textwidth}{c X}
\toprule
\textbf{1d6} & \textbf{Hallazgos en el Asfalto y Bolsillos de la Margen} \\
\midrule
\textbf{1} & Una petaca de aguardiente casero con la inscripción «Turno de Noche». \\
\textbf{2} & Una tarjeta perforada con códigos de anulación de una esclusa de ventilación. \\
\textbf{3} & Tres viales de suero antitóxico caducado pero todavía útil. \\
\textbf{4} & Una fotografía en blanco y negro arrugada de la Ría limpia hace medio siglo. \\
\textbf{5} & Un detonador artesanal de plasma fabricado con condensadores de CRT. \\
\textbf{6} & Una chapa de afiliación sindical clandestina con el símbolo de la forja. \\
\bottomrule
\end{tabularx}
\caption{Tabla de Chatarra y Objetos de Valor.}
\end{table}

\begin{table}[htbp]
\centering
\small
\begin{tabularx}{\textwidth}{c X}
\toprule
\textbf{1d6} & \textbf{Rumores en las Tabernas y Frecuencias Clandestinas} \\
\midrule
\textbf{1} & Se dice que en la planta -4 han encontrado a un dron Tártalo pintado de rojo. \\
\textbf{2} & La ración de pan sintético bajará un 15\% a partir de la próxima semana. \\
\textbf{3} & Una furgoneta azul sin matrícula fue vista merodeando la escuela de Sestao. \\
\textbf{4} & Dicen que la señal de radio pirata emite canciones prohibidas a medianoche. \\
\textbf{5} & Los auditores de la Cúpula han duplicado las cuotas de Burdin-Urdin del mes. \\
\textbf{6} & Un grupo de soldadores del Alto Horno ha jurado no entregar a sus hijos. \\
\bottomrule
\end{tabularx}
\caption{Tabla de Rumores y Tensión Callejera.}
\end{table}

\newpage

\section{Hoja de Referencia Rápida de Reglas (Cheat Sheet)}

\begin{CajaRegla}{Mecánica Fundamental ZINA}
\textbf{Tirada básica:} Lanza \textbf{2d6 + Emoción relevante} (+1d4 de Auzolan si te asisten).
\begin{itemize}
    \item \textbf{10 o más (Éxito Rotundo):} Logras lo que pretendías de forma limpia y obtienes una ventaja posicional o táctica añadida.
    \item \textbf{7 a 9 (Éxito Parcial o con Coste):} Consigues tu objetivo, pero sufres una complicación, daño menor, gasto de recursos o marcas Zama.
    \item \textbf{6 o menos (Fallo Dramático):} No consigues tu objetivo de forma limpia o empeora gravemente la situación. La DJ realiza un Movimiento Duro y marcas 1 casilla de Zama.
\end{itemize}
\end{CajaRegla}

\begin{CajaRegla}{Los Seis Movimientos Comunes}
\begin{enumerate}
    \item \textbf{Afrontar la Presión (Fatiga):} Resistir dolor, sostener pesos, superar turnos dobles.
    \item \textbf{Golpear con Furia (Rencor):} Violencia directa, combates cuerpo a cuerpo, amenazas.
    \item \textbf{Forzar la Máquina (Esperanza):} Hacking analógico CRT, sobrecargas, puentear el Velo.
    \item \textbf{Otear el Asfalto (Arraigo):} Leer la calle, detectar trampas corporativas, encontrar contactos.
    \item \textbf{Activar el Auzolan (Arraigo / Reserva):} Pedir auxilio a la cuadrilla, añadir +1d4 de apoyo.
    \item \textbf{Resistir la Fisura (Esperanza / Vínculo):} Contener el pánico y evitar marcar Zama.
\end{enumerate}
\end{CajaRegla}

\begin{CajaRegla}{Gestión de Zama y Quiebre Emocional}
\begin{itemize}
    \item Cada personaje tiene \textbf{5 casillas de Zama}.
    \item Al marcar la 5.ª casilla, entras en \textbf{Punto de Quiebre}: debes elegir entre \textit{Estallido de Furia} (violencia incontrolable hacia la autoridad más cercana), \textit{Parálisis Catatónica} (incapaz de actuar durante una escena) o \textit{Ruptura de Vínculo} (destruir la relación de confianza con un compañero).
    \item La Zama se recupera descansando en un refugio seguro, compartiendo comida y bebida con la cuadrilla o cumpliendo una parte de tu Juramento (\textit{Zina}).
\end{itemize}
\end{CajaRegla}
